\section{Data Analysis}

(HOW THE ANALYSIS WILL BE CARRIED OUT AND WHAT BASIC TESTS WILL BE USED) AND PRELIMINARY RESULTS OR PILOT TESTS (IF AVAILABLE)

This section describes whether a pilot test is planned, how data collectors will be trained and how the data will be coded. Additionally, it should include an explanation of how access to the relevant information will be ensured.

If the project involves experimental designs, the procedures must clearly explain the nature of the treatment(s), the levels of the independent variable, the number of groups involved, the method for assigning participants or cases to groups, your role as a researcher in the experiment, the time span between pre-test, treatment, and post-test for each group, the way in which the groups will be monitored, and the location where the study will be conducted.

Regarding data analysis, Mertens\cite{mertens2015} includes the following:

\begin{itemize}
	\item \textit{Plan for processing the data:} coding, analytical software, and statistical tests to be used (for each hypothesis or variable; if no hypotheses were stated, then for each research question and variable).
	\item The way you will link the analysis or analyses to the research problem and hypotheses (it is not enough to mention which tests will be applied; you must establish the connection and the type of results expected).
\end{itemize}

Dahlberg, Wittink, and Gallo\cite{dahlberg2010} recommend adding, at the end of the methods section, a paragraph that explains the strengths and challenges of your research. This should include a description of alternative methods that could be considered and explain why the chosen method represents the best approach to the research problem. Alternatively, it should address how potential threats to the project and its feasibility will be managed. This paragraph demonstrates to reviewers and readers that you have anticipated problems, obstacles, and criticisms.

Additionally, you must honestly acknowledge the limitations of the study—something highly valued by all researchers.
